\appendix

\noindent\emph{Organization: design rationale (A$^*$), confound analysis (A--B), model diagnostics (C), independence tests (D), gate validation (E--F), equivalence testing (G), statistical details (H--K), sensitivity analyses (L--N), artifact controls (O), implementation (P), multiplicity (Q), cross-axis contamination (R), FK length-control (R$'$), human annotation agreement (S), cross-benchmark replication (T), Claude cross-benchmark (U), rewriting prompts (V), non-tie concordance (W).}

\section{Design Rationale}
\label{app:design_rationale}

Three variables govern the observed judge score $Y$: the style attribute $S$ (e.g., formality), the original content quality $Q$, and the judge's internal scoring function.
Both $S$ and $Q$ causally influence $Y$; the goal of CSD is to approximate the effect $S \to Y$ while blocking confounding through $Q$.
Style-controlled rewriting generates a text pair that differs along a single style axis, but any rewriting operation risks inadvertently altering content quality ($R \to Q$).
The NLI gate (\S\ref{sec:nli_gate}) blocks this path, yielding a controlled estimate of the $S \to Y$ effect.
Our approach provides interventional rather than formal causal evidence: we do not claim that the DAG in Figure~\ref{fig:dag} is a complete structural causal model, nor do we invoke SUTVA or consistency assumptions from the potential outcomes framework.
``Controlled'' in this context refers to experimentally controlled rewriting conditions, not formal causal identification in the potential-outcomes sense.
Our identification rests on three operational assumptions: (i)~the NLI gate approximates content preservation, blocking the $R \to Q$ path; (ii)~position swapping controls for presentation-order effects; and (iii)~the style rewrite introduces no systematic differences beyond the target axis.

\section{Direction Asymmetry and Quality Confounding}
\label{app:direction}

The formality axis exhibits a pronounced direction asymmetry: the judge selects the formal variant in 78.8\% of formal-is-rewrite pairs ($p = 0.001$) but only 54.5\% of casual-is-rewrite pairs ($p = 0.83$).
This asymmetry suggests that original quality $Q$ and formality $S$ are partially confounded, as formally written texts may be inherently perceived as higher quality.
To quantify this, we fit a GEE logistic model (exchangeable correlation) with both formality and original-preference as predictors.
After controlling for original preference, the formality effect remains significant ($\text{OR} = 2.58$, $p = 0.002$), and the original-preference effect is also significant ($\text{OR} = 2.32$, $p = 0.006$).
The formality preference thus survives adjustment for original-preference confounding, though the significant original-preference effect indicates that quality--style confounding is present and non-negligible.

At $n = 88$, the Qwen3-32B formality Style OR ($4.20$) exceeded the GPT-4o estimate ($2.58$), but a scaled replication ($n = 355$) corrects this to $2.15$ ($p < 0.0001$), comparable to GPT-4o and consistent with small-sample upward bias rather than evaluator-producer inflation of the style effect.
The originality confound remains stronger ($4.99$ vs.\ $2.32$), consistent with evaluator-producer bias inflating the originality component (Limitations).
The directional agreement between judges (both favor formal text after controlling for originality) supports the existence of formality bias, but the effect magnitude from Qwen3-32B should not be directly compared with GPT-4o.

The NLI gate blocks \emph{content} changes but cannot block quality perception shifts that are themselves mediated by style, a fundamental identifiability limit of any text-rewriting approach.

\section{Consistency Filter Robustness}
\label{sec:consistency}

The position-consistency filter retains 55 of 98 valid pairs (56.1\%); consistent and inconsistent groups show no significant differences across four covariates (NLI score $p = 0.824$, BERTScore $p = 0.12$, text length $p = 0.137$, length diff $p = 0.715$), indicating selection on judge decisiveness rather than text properties.
The BT model using all 173 comparisons including inconsistent pairs yields 63.0\% ($p = 0.0004$), consistent with the filtered sign test (73.9\%, $p = 0.006$), confirming no systematic bias from the filter.
Chi-squared tests confirm that NLI gate pass/reject status is independent of source category ($\chi^2 = 1.30$, $p = .73$), original formality ($p = .37$), and gold label ($p = 1.0$), indicating no systematic selection bias along these dimensions.

Verbosity gate-passed pairs are significantly longer than rejected pairs (token count KS $p = 0.003$; sentence length MWU $p = 0.006$).
This reflects the gate's content-preservation mandate: verbosity transformations on shorter texts alter content ratios more substantially, causing correct NLI rejection.
Length-controlled GEE analysis (\S\ref{sec:direction}; Style OR change ${<}2\%$) confirms this selection does not confound the verbosity effect.

\section{GEE Model Diagnostics}
\label{app:gee_diagnostics}

\paragraph{Correlation structure robustness.}
Table~\ref{tab:gee_corr} compares exchangeable and independence working correlation structures across all nine judge$\times$axis cells.
The two structures yield virtually identical results: $|\Delta\text{QIC}| < 0.001$ and maximum coefficient change ${<}0.02$ in every cell, with estimated within-pair correlations $\hat{\alpha} \in [-0.25, +0.11]$.
All GEE results reported in the main text are therefore robust to correlation structure specification.

\begin{table}[h]
\centering
\small
\caption{GEE correlation structure comparison across nine judge$\times$axis cells.}
\label{tab:gee_corr}
\resizebox{\columnwidth}{!}{\begin{tabular}{@{}lcccc@{}}
\toprule
Config & Exch.\ QIC & Indep.\ QIC & $|\Delta\text{QIC}|$ & Max$|\Delta\text{Coef}|$ \\
\midrule
GPT-4o form.\ & --- & --- & ${<}0.001$ & ${<}0.02$ \\
Qwen3 form.\ & --- & --- & ${<}0.001$ & ${<}0.02$ \\
Llama-70B form.\ & --- & --- & ${<}0.001$ & ${<}0.02$ \\
GPT-4o verb.\ & --- & --- & ${<}0.001$ & ${<}0.02$ \\
Qwen3 verb.\ & --- & --- & ${<}0.001$ & ${<}0.02$ \\
Llama-70B verb.\ & --- & --- & ${<}0.001$ & ${<}0.02$ \\
GPT-4o reg.\ & --- & --- & ${<}0.001$ & ${<}0.02$ \\
Qwen3 reg.\ & --- & --- & ${<}0.001$ & ${<}0.02$ \\
Llama-70B reg.\ & --- & --- & ${<}0.001$ & ${<}0.02$ \\
\bottomrule
\multicolumn{5}{l}{\footnotesize Absolute QIC values omitted; only $|\Delta\text{QIC}|$ is reported.}
\end{tabular}}
\end{table}

\paragraph{Interaction test.}
Adding a Style$\times$Originality interaction term to the additive GEE specification yields no significant interactions in any of the nine judge$\times$axis cells (all $p > 0.13$, range $0.14$--$0.96$), confirming that the additive model is the correct specification and that regime classifications are robust to model specification.
Table~\ref{tab:gee-interaction} reports interaction terms; none are significant, confirming the additive specification.

\begin{table}[h]
\centering
\small
\caption{GEE Style $\times$ Originality interaction terms. None are significant ($p > 0.13$), confirming the additive model specification.}
\label{tab:gee-interaction}
\resizebox{\columnwidth}{!}{\begin{tabular}{@{}llccc@{}}
\toprule
Judge & Axis & Interaction OR & 95\% CI & $p$ \\
\midrule
GPT-4o & Formality & 0.661 & [0.18, 2.38] & .527 \\
GPT-4o & Verbosity & 1.269 & [0.37, 4.34] & .705 \\
GPT-4o & Register & 2.601 & [0.74, 9.15] & .136 \\
Qwen3-32B & Formality & 0.518 & [0.14, 1.89] & .320 \\
Qwen3-32B & Verbosity & 0.450 & [0.13, 1.62] & .221 \\
Qwen3-32B & Register & 0.555 & [0.15, 2.08] & .383 \\
Llama-70B & Formality & 0.968 & [0.25, 3.78] & .963 \\
Llama-70B & Verbosity & 1.103 & [0.34, 3.58] & .870 \\
\bottomrule
\end{tabular}}
\end{table}

\paragraph{95\% confidence intervals.}
Table~\ref{tab:gee_ci} reports 95\% CIs for all GEE odds ratios in Table~\ref{tab:gee}.

\begin{table}[h]
\centering
\small
\caption{95\% CIs for GEE conditional decomposition (Table~\ref{tab:gee}).}
\label{tab:gee_ci}
\begin{tabular}{@{}lcc@{}}
\toprule
Config & Style OR [CI] & Orig OR [CI] \\
\midrule
Qwen3 verb.\ & 13.10 [5.62, 30.55] & 0.49 [0.21, 1.14] \\
GPT-4o verb.\ & 11.34 [4.93, 26.10] & 0.45 [0.20, 1.04] \\
Llama-70B verb.\ & 8.21 [3.62, 18.65] & 0.38 [0.17, 0.86] \\
Gemini verb.\ & 8.98 [3.68, 21.94] & 0.46 [0.20, 1.04] \\
GPT-4o form.\ & 2.58 [1.41, 4.73] & 2.32 [1.27, 4.26] \\
Qwen3 form.\ & 4.20 [1.99, 8.87] & 4.70 [2.22, 9.96] \\
Qwen3 form.$^a$ & 2.15 [1.46, 3.16] & 4.99 [3.40, 7.34] \\
Llama-70B form.\ & 4.81 [2.48, 9.33] & 3.12 [1.61, 6.04] \\
Gemini form.\ & 1.64 [0.91, 2.98] & 2.51 [1.39, 4.54] \\
GPT-4o reg.\ & 1.88 [1.08, 3.28] & 2.21 [1.26, 3.86] \\
Qwen3 reg.\ & 1.91 [1.01, 3.58] & 5.72 [3.02, 10.82] \\
Llama-70B reg.\ & 3.33 [1.73, 6.42] & 3.07 [1.60, 5.92] \\
\bottomrule
\multicolumn{3}{l}{\footnotesize $^a$ $n = 355$; formality rows: $n = 88$; verbosity rows: $n = 86$.}
\end{tabular}
\end{table}

\section{Rewriter Independence}
\label{app:rewriter_indep}

Evaluator-producer overlap is limited to Qwen3-32B (rewriter $=$ judge); the independence test below bounds this to a 22\% Style OR reduction. The other three target judges (GPT-4o, Llama-70B, Gemini) use a different rewriter and are not subject to this confound.

To test whether the Blindspot depends on the rewriter's identity, we replace Qwen3-32B with GPT-4o as rewriter while keeping Qwen3-32B as judge.

\paragraph{Verbosity (Table~\ref{tab:rewriter_indep}).}
The Style OR decreases by $22\%$ ($13.10 \to 10.17$), bounding evaluator-producer inflation of the style effect; the pure-style pattern (Originality n.s.) and Blindspot (BT significant, Likert n.s.) both persist.

\begin{table}[h]
\centering
\small
\caption{Rewriter independence test: verbosity axis, Qwen3-32B judge. GPT-4o rewriter eliminates evaluator-producer overlap while preserving the Blindspot.}
\label{tab:rewriter_indep}
\begin{tabular}{lcc}
\toprule
 & GPT-4o Rewriter & Qwen3 Rewriter \\
\midrule
NLI pass rate & 81/100 (81\%) & 86/100 (86\%) \\
BT $P$ [95\% CI] & .735 [.654, .809] & .773 [.703, .843] \\
Likert ATE ($p$) & .000 (.838) & ${\sim}$0 (.179) \\
Blindspot & Yes & Yes \\
GEE Style OR ($p$) & 10.17 (${<}.0001$) & 13.10 (${<}.0001$) \\
GEE Orig OR ($p$) & 1.61 (n.s.) & 0.49 (n.s.) \\
\bottomrule
\end{tabular}
\end{table}

\paragraph{Formality (Table~\ref{tab:rewriter_indep_form}).}
GPT-4o rewriter produces formality-manipulated pairs yielding BT $64.4\%$ $[56.1, 72.8]$, statistically indistinguishable from the Qwen3-32B baseline ($66.5\%$ $[58.5, 73.9]$; $\Delta = -2.1$pp).
The Likert Blindspot replicates with the independent rewriter ($p = 0.184$, n.s.).

\begin{table}[h]
\centering
\small
\caption{Rewriter independence test: formality axis, Qwen3-32B judge.}
\label{tab:rewriter_indep_form}
\begin{tabular}{lcc}
\toprule
 & GPT-4o Rewriter & Qwen3 Rewriter \\
\midrule
NLI pass rate & 90/100 (90\%) & 88/100 (88\%) \\
BT $P$ [95\% CI] & .644 [.561, .728] & .665 [.585, .739] \\
Likert $d$ ($p$) & $-$.081 (.184) & .011 (.908) \\
Blindspot & Yes & Yes \\
\bottomrule
\end{tabular}
\end{table}

\section{NLI Threshold Sensitivity Sweep}
\label{app:nli_sweep}

Table~\ref{tab:nli_sweep} reports BT win probabilities across 11 NLI entailment thresholds for all four judge$\times$axis configurations.
All 44 combinations are significant ($p < 0.05$), and within each configuration BT varies by at most 3pp, ruling out threshold selection bias.
The stability reflects the bimodal NLI score distribution: most pairs score above 0.95 or below 0.30, so thresholds between 0.50 and 0.95 admit nearly identical pair sets.

\begin{table}[h]
\centering
\small
\caption{BT win probability across NLI thresholds. All 44 cells significant ($p < 0.05$). $n$: pairs passing threshold.}
\label{tab:nli_sweep}
\begin{tabular}{@{}rcccc@{}}
\toprule
Thresh & \multicolumn{2}{c}{Formality BT} & \multicolumn{2}{c}{Verbosity BT} \\
\cmidrule(lr){2-3} \cmidrule(lr){4-5}
 & Qwen3 & GPT-4o & Qwen3 & GPT-4o \\
\midrule
.50 & .648 & .615 & .770 & .768 \\
.55 & .648 & .615 & .770 & .768 \\
.60 & .656 & .621 & .770 & .768 \\
.65 & .656 & .621 & .767 & .766 \\
.70 & .657 & .623 & .767 & .766 \\
.75 & .657 & .623 & .773 & .760 \\
.80 & .657 & .623 & .773 & .760 \\
.85 & .657 & .623 & .773 & .760 \\
.90 & .665 & .630 & .773 & .760 \\
.95 & .653 & .629 & .765 & .752 \\
.98 & .649 & .642 & .771 & .769 \\
\bottomrule
\end{tabular}
\end{table}

\section{Direction-Aware Gate Validation}
\label{app:gate_validation}

For verbosity, the direction-aware NLI gate admits 21 additional pairs (32\%) compared to strict bidirectional verification.
The two groups show no significant difference in BT win probability (78.5\% bidirectional vs.\ 73.8\% DA-only; permutation $p = .663$; $\Delta = {-}1.1$pp, 95\% CI $[{-}5.6, {+}3.0]$) or Likert $d$ ($0.038$ vs.\ $0.029$).
Both groups exhibit the Blindspot pattern (BT significant, Likert NS), confirming that direction-aware gating recovers valid pairs without inflating the effect or introducing selection bias.
Cross-judge validation on all 21 direction-aware pairs confirms this: Fisher exact tests comparing BT estimates between direction-aware and bidirectional-pass subsets are non-significant for all five judges ($p > 0.66$), indicating that the DA-gate does not differentially affect any judge's preference signal.
For formality, strict bidirectional verification is appropriate: a relaxed (direction-aware) gate would admit 9 additional pairs with reversed preference (BT 33.4\%), indicating content preservation failure rather than genuine style preference.
This asymmetry empirically justifies the axis-dependent gating strategy described in Section~\ref{sec:nli_gate}.

\section{Claude Sonnet 4 Equivalence Testing}
\label{app:claude_tost}

Claude Sonnet~4 exhibits benchmark-dependent sensitivity ($n = 88$; LLMBar BT 56.2\%/54.1\%, $p = 0.091$/$0.284$).
Formal equivalence testing \citep[TOST;][]{lakens2017equivalence} yields margin-dependent conclusions (Table~\ref{tab:tost_sensitivity}): per-axis equivalence is confirmed at $\Delta = 15$pp ($p = .010$), corresponding to a BT threshold of 65\%.
At $\Delta = 10$pp, per-axis TOST is not significant ($p = .156$, power $\approx 42\%$).
Pooling both axes ($n = 174$, BT $= 55.17\%$) achieves equivalence at $\Delta = 12$pp ($p = .035$) but not at $\Delta = 10$pp ($p = .100$), tightening the bound on Claude's style sensitivity.
We adopt $\Delta = 12$pp as it corresponds to a small effect (Cohen's $h = 0.24$), below which style sensitivity is unlikely to materially affect downstream RLHF reward modeling; a 12pp BT shift means ${\sim}12\%$ of preference pairs change direction. This is also the smallest margin at which pooled significance is achieved.
We report all tested margins for transparency.

\begin{table}[h]
\centering
\small
\caption{TOST equivalence sensitivity for Claude. Per-axis: formality only (BT $= 0.562$, $n = 88$). Pooled: both axes combined ($n = 174$).}
\label{tab:tost_sensitivity}
\begin{tabular}{ccl}
\toprule
$\Delta$ (pp) & $p$ & Result \\
\midrule
5 & .626 & Not equivalent (per-axis) \\
10 & .156 & Not equivalent (per-axis) \\
10 & .100 & Not equivalent (pooled) \\
12 & .035 & Equivalent (pooled) \\
15 & .010 & Equivalent (per-axis) \\
20 & .0001 & Equivalent (per-axis) \\
\bottomrule
\end{tabular}
\end{table}

\paragraph{Bayesian equivalence.}
We complement frequentist TOST with Bayesian equivalence testing to quantify evidence strength for the null hypothesis.
Table~\ref{tab:bayesian_equiv} reports BF$_{01}$ (evidence for null over alternative) under three half-normal prior scales.

\begin{table}[h]
\centering
\small
\caption{Bayesian equivalence (BF$_{01}$) for Claude Sonnet~4 Likert ATE. BF$_{01} > 3$: moderate evidence for null; $1$--$3$: anecdotal; $< 1$: favors alternative.}
\label{tab:bayesian_equiv}
\begin{tabular}{@{}lccc@{}}
\toprule
Data & $\sigma = 0.3$ & $\sigma = 0.5$ & $\sigma = 1.0$ \\
\midrule
Pooled ($n = 348$) & 0.57 & 0.80 & 1.48 \\
Verbosity & 1.40 & 2.03 & 3.78 \\
Formality & 0.74 & 0.97 & 1.73 \\
\bottomrule
\end{tabular}
\end{table}

Pooled evidence is inconclusive (BF$_{01} < 3$ across all priors), consistent with Claude exhibiting a small (${\sim}5$pp) rather than zero style preference on LLMBar.
Verbosity alone approaches moderate evidence for the null (BF$_{01} = 2.03$--$3.78$); formality pulls down the pooled estimate (BF$_{01} = 0.74$--$1.73$).

\section{Bradley-Terry Goodness of Fit and Stratified Analysis}
\label{app:bt_gof}

Bradley-Terry goodness-of-fit tests confirm adequate model fit for both axes (formality: deviance $= 52.80$, $\text{df} = 45$, $p = 0.198$; verbosity: $p = 0.905$).
Stratified analysis of consistent pairs by original-text direction reveals no preference heterogeneity for verbosity (Fisher $p = 0.453$) but significant direction asymmetry for formality ($88.9\%$ vs.\ $52.6\%$, Fisher $p = 0.015$), independently confirming the originality confound structure identified by GEE (Table~\ref{tab:gee}).

To assess robustness to originality confounding, we restrict to pairs where the judge assigned equal Likert scores to original and counterfactual texts (eliminating detectable quality differences).
Style preference persists: formality BT $= 67.3\%$ ($n = 55$, $\Delta = {+}0.8$pp from full sample, 95\% CI $[{-}11.0, {+}12.6]$); verbosity BT $= 76.0\%$ ($n = 50$, $\Delta = {-}1.3$pp, 95\% CI $[{-}13.3, {+}10.3]$).
In the anti-originality subset (pairs where the original never wins), effects strengthen (formality $75.0\%$, verbosity $86.1\%$), confirming that style preference is not an artifact of originality confounding.

\section{Bradley-Terry Confidence Intervals}
\label{app:bt_ci}

Table~\ref{tab:bt_ci} reports bootstrap 95\% confidence intervals for all Bradley-Terry estimates in Table~\ref{tab:main_results}.

\begin{table}[h]
\centering
\small
\caption{Bootstrap 95\% CIs for BT win probabilities (10{,}000 pair-level resamples).}
\label{tab:bt_ci}
\begin{tabular}{lllc}
\toprule
Bench & Judge & Axis & BT $P$ [95\% CI] \\
\midrule
LLMBar & GPT-4o & Form.\ & .630 [.557, .703] \\
 & GPT-4o & Verb.\ & .760 [.667, .853] \\
 & GPT-4o & Reg.\ & .570 [.468, .673] \\
\cmidrule{2-4}
 & Qwen3 & Form.\ & .665 [.585, .739] \\
 & Qwen3 & Verb.\ & .773 [.703, .843] \\
 & Qwen3 & Reg.\ & .553 [.479, .628] \\
\cmidrule{2-4}
 & Llama-70B & Form.\ & .695 [.619, .765] \\
 & Llama-70B & Verb.\ & .735 [.655, .809] \\
 & Llama-70B & Reg.\ & .635 [.555, .712] \\
\cmidrule{2-4}
 & Claude & Form.\ & .562 [.489, .636] \\
 & Claude & Verb.\ & .541 [.459, .616] \\
\midrule
AEval & GPT-4o & Form.\ & .727 [.663, .785] \\
 & Qwen3 & Form.\ & .657 [.581, .733] \\
 & Qwen3 & Verb.\ & .841 [.774, .902] \\
\midrule
MTB & GPT-4o & Form.\ & .692 [.623, .753] \\
 & Qwen3 & Form.\ & .712 [.630, .788] \\
\bottomrule
\end{tabular}
\end{table}

\section{Bootstrap Attenuation Analysis}
\label{app:bootstrap_attenuation}

To distinguish measurement attenuation from insufficient statistical power, we subsample the $n = 355$ formality replication (Qwen3-32B) at seven sample sizes and compute Likert Cohen's $d$ with bootstrap 95\% CIs at each (Table~\ref{tab:attenuation}).
Under a power limitation, $d$ would remain constant while only $p$-values improve with larger $n$; under measurement attenuation, $d$ itself converges to zero.
The data support attenuation: the bootstrap median $d = {-}0.017$ across all sample sizes (the point estimate at $n = 355$ is $d = {-}0.025$; Table~\ref{tab:main_results}), with CIs narrowing from ${\pm}0.3$ to ${\pm}0.1$ but never excluding zero.
The same data yield BT $= 67.5\%$ ($p < 0.0001$), confirming that the pairwise signal is robust.
Verbosity shows a similar pattern ($d \approx 0.15$--$0.19$, stable, CIs always containing zero; BT $= 76$--$77\%$, $p < 0.0001$).
Within the agnostic framework of \S\ref{sec:likert_pairwise}, this evidence favors the measurement attenuation interpretation but does not exclude construct divergence.

\begin{table}[h]
\centering
\small
\caption{Bootstrap attenuation analysis: formality $\times$ Qwen3-32B. Bootstrap median Likert $d$ remains flat near zero across sample sizes while BT is stable at 67.5\%. Point estimate at $n = 355$: $d = {-}0.025$ (Table~\ref{tab:main_results}).}
\label{tab:attenuation}
\begin{tabular}{@{}rccc@{}}
\toprule
$n$ & Likert $d$ & 95\% CI & Contains 0? \\
\midrule
50 & $-$0.017 & [$-$0.308, 0.262] & Yes \\
100 & $-$0.015 & [$-$0.219, 0.176] & Yes \\
150 & $-$0.019 & [$-$0.179, 0.145] & Yes \\
200 & $-$0.016 & [$-$0.157, 0.119] & Yes \\
250 & $-$0.018 & [$-$0.144, 0.107] & Yes \\
300 & $-$0.018 & [$-$0.131, 0.095] & Yes \\
355 & $-$0.017 & [$-$0.123, 0.087] & Yes \\
\bottomrule
\end{tabular}
\end{table}

\section{Formal Power Analysis}
\label{app:power_analysis}

Table~\ref{tab:power} reports the sample size required for Likert ATE to achieve 80\% power at observed effect sizes across all six judge$\times$axis cells.
Required sample sizes range from $336$ (Qwen3-32B verbosity, $d = 0.153$) to $27{,}159$ (Qwen3-32B formality, $d = {-}0.017$), representing $4$--$206\times$ the current $n$.
For five of six cells, the required $n$ exceeds $1{,}200$, far beyond typical evaluation budgets, converging with the bootstrap attenuation curve (Appendix~\ref{app:bootstrap_attenuation}) to confirm the Blindspot as a structural measurement limitation.

\begin{table}[h]
\centering
\small
\caption{Required Likert sample size for 80\% power at observed $d$. BT detects all six effects at $p < 0.001$ with $n \leq 100$.}
\label{tab:power}
\begin{tabular}{@{}llcrc@{}}
\toprule
Judge & Axis & Observed $d$ & Required $n$ & Ratio \\
\midrule
GPT-4o & Form.\ & 0.073 & 1{,}473 & 16.7$\times$ \\
GPT-4o & Verb.\ & 0.080 & 1{,}227 & 14.3$\times$ \\
Qwen3 & Form.\ & $-$0.017 & 27{,}159 & 76.5$\times$ \\
Qwen3 & Verb.\ & 0.153 & 336 & 3.9$\times$ \\
Llama-70B & Form.\ & 0.021 & 18{,}157 & 206.3$\times$ \\
Llama-70B & Verb.\ & 0.034 & 6{,}737 & 78.3$\times$ \\
\bottomrule
\end{tabular}
\end{table}

\section{Formality Sensitivity Hierarchy}
\label{app:sensitivity}

Table~\ref{tab:sensitivity} applies seven measurement methods to the same formality style-transformed data (GPT-4o, LLMBar, $n = 88$).
Only pairwise methods reach significance; all Likert-based methods fail regardless of scaling assumption.

\begin{table}[h]
\centering
\small
\caption{Sensitivity hierarchy: seven measurement methods applied to the same formality data. Only pairwise methods (top two rows) reach significance.}
\label{tab:sensitivity}
\resizebox{\columnwidth}{!}{\begin{tabular}{@{}lccc@{}}
\toprule
Method & Statistic & $p$ & Data \\
\midrule
BT model & 63.0\% $[55.7, 70.3]$ & $0.0004$ & All 173 comp. \\
Sign test (filtered) & 73.9\% (34/46) & 0.006 & Consistent \\
\midrule
Likert ATE ($t$-test) & $d = 0.102$ & 0.332 & NLI$\geq$0.90 \\
Binarized Likert & 60.6\% & 0.296 & Non-tied \\
Ordinal reg.\ (prop.\ odds) & $\beta = {-}0.079$ & 0.769 & NLI$\geq$0.90 \\
GEE (exchangeable) & $\beta = {-}0.084$ & 0.587 & NLI$\geq$0.90 \\
Wilcoxon signed-rank & $W = 274$ & 0.905 & NLI$\geq$0.90 \\
\bottomrule
\end{tabular}}
\end{table}

\section{Verbosity Sensitivity Hierarchy}
\label{app:verb_sensitivity}

Table~\ref{tab:verb_sensitivity} replicates the seven-method sensitivity analysis (Table~\ref{tab:sensitivity}) on the verbosity axis ($n = 86$ NLI-filtered pairs, Qwen3-32B judge, LLMBar).
The attenuation pattern is more extreme than formality: all six Likert-based methods remain non-significant (lowest $p = 0.213$) despite a stronger pairwise signal (BT 77.3\% vs.\ 66.5\% for formality), and Likert effect sizes ($d = 0.036$--$0.056$) are near zero.

\begin{table}[h]
\centering
\small
\caption{Sensitivity hierarchy for verbosity ($n = 86$, Qwen3-32B, LLMBar). All Likert methods fail; only BT reaches significance.}
\label{tab:verb_sensitivity}
\begin{tabular}{lccc}
\toprule
Method & Statistic & $p$ & Sig? \\
\midrule
BT model & 77.3\% $[70.3, 84.3]$ & ${<}0.001$ & Yes \\
\midrule
Likert ATE & $d = 0.056$ & 0.691 & No \\
Binarized $\chi^2$ & OR $= 2.16$ & 0.213 & No \\
Ordinal reg.\ & $\beta = 0.159$ & 0.575 & No \\
Paired $t$-test & $d = 0.056$ & 0.607 & No \\
Mann-Whitney $U$ & $r = {-}0.046$ & 0.578 & No \\
Cohen's $d$ & $d = 0.036$ & --- & No \\
\bottomrule
\multicolumn{4}{p{0.9\columnwidth}}{\footnotesize Effect sizes use style-aligned scoring (verbose $-$ concise), isolating style preference direction. Table~\ref{tab:main_results} reports RATE-compatible ATE (counterfactual $-$ original), which conflates style and originality effects.}
\end{tabular}
\end{table}

\subsection{Scaled Replication ($n = 330$)}
\label{app:verb_n355}

Table~\ref{tab:verb_n355} reports verbosity results at $n = 330$ (from 355 generated pairs, 330 passing NLI gate at $\tau = 0.90$), replicating across four judges (GPT-4o was not re-scored at this scale).
The Blindspot partially resolves for judges with strong style effects: Qwen3-32B ($d = .149$, $p < .001$) and Llama-70B ($d = .160$, $p < .001$) now reach Likert significance, consistent with the power curve (Figure~\ref{fig:power_curve}) predicting ${\sim}40\%$ power at $n = 330$ for $d \approx 0.15$.
However, two judges maintain the Blindspot: Gemini-2.5-Flash ($d = .087$, borderline) and Claude Sonnet~4 ($d = .032$, NS).
The magnitude discrepancy persists even where Likert reaches significance: BT detects 73--76\% win rates while Likert yields $d \approx 0.15$, confirming the attenuation interpretation.

GEE conditional decomposition at $n = 330$ confirms substantial shrinkage from $n = 88$ estimates, consistent with more extreme initial estimates producing larger corrections \citep{nemes2009bias}: Qwen3 $13.10 \to 3.83$ ($-71\%$), Llama $8.21 \to 4.01$ ($-51\%$), Gemini $8.98 \to 2.52$ ($-72\%$).
Despite shrinkage, verbosity remains style-dominated for these three judges (Style OR $2.5$--$4.0$, all $p < .001$; Orig OR $0.52$--$0.69$).
Claude shows a non-significant style effect at $n = 330$ (Style OR $= 1.22$, $p = .134$), with its elevated BT reflecting originality rather than style preference (Orig OR $= 1.69$).
As predicted by the power analysis (Figure~\ref{fig:power_curve}), Likert begins detecting the strongest effects at $n \approx 330$ (Qwen3 $d = .15$, $p < .001$; Llama $d = .16$, $p < .001$), while two judges maintain the Blindspot (Claude $d = .03$, NS; Gemini borderline $d = .09$, $p = .02$ by Wilcoxon but $p = .11$ by $t$-test).
The efficiency gap (Likert requiring $3$--$4\times$ the sample size pairwise needs), rather than absolute failure, is the practical concern.

\begin{table}[h]
\centering
\small
\caption{Verbosity scaled replication ($n = 330$, four judges; GPT-4o not re-scored). Blindspot partially resolves for strong effects (Qwen3, Llama-70B) but persists for weaker-effect judges (Gemini, Claude).}
\label{tab:verb_n355}
\resizebox{\columnwidth}{!}{\begin{tabular}{@{}lcccc@{}}
\toprule
Judge & BT $P(\text{V})$ & [95\% CI] & Likert $d$ ($p$) & Blindspot? \\
\midrule
Qwen3-32B & .733$^{***}$ & [.696, .770] & .149 (.001) & No \\
Llama-70B & .759$^{***}$ & [.722, .793] & .160 ({<}.001) & No \\
Gemini-2.5 & .678$^{***}$ & [.638, .717] & .087 (.021) & $\dagger$ \\
Claude S4 & .608$^{***}$ & [.567, .649] & .032 (.317) & $\dagger$ \\
\bottomrule
\end{tabular}}
\par\vspace{2pt}
{\footnotesize $^{***}$\,$p < 0.001$. $\dagger$ marks Blindspot cases (BT significant, Likert $d$ non-significant at $\alpha = 0.05$ by Wilcoxon). Likert $p$ values from Wilcoxon signed-rank test on non-zero differences.}
\end{table}

\section{100-Point Continuous Likert Scale}
\label{app:likert100}

To test whether the Blindspot reflects integer-scale granularity rather than a paradigm difference between absolute and relative evaluation, we replicate the formality experiment using a 100-point continuous scale ($1$--$100$) with three judges (Table~\ref{tab:likert100}).
All three judges remain non-significant (all $p > 0.15$).
Tie rates decrease substantially with finer granularity (1--10 scale: 58--73\%; 1--100 scale: 40--66\%), yet ATE remains non-significant, confirming the gap reflects measurement mode rather than scale granularity \citep{preston2000scale}.
Judges differ markedly in effective granularity usage: GPT-4o clusters scores into ${\sim}5$ distinct bins (modes at 70, 75, 80), Qwen3-32B uses ${\sim}20$ distinct values spanning $45$--$95$, and Llama-70B exploits $50{+}$ unique scores with near-uniform coverage, yet none resolves the Blindspot.
Score distributions are unimodal and right-skewed for all three judges, with medians in the $72$--$80$ range; the score \emph{difference} distributions (original minus rewrite) are symmetric around zero, confirming no net directional shift under Likert scoring.

\begin{table}[h]
\centering
\small
\caption{Formality Likert ATE: 1--10 vs.\ 1--100 scale. Finer granularity reduces tie rates but does not resolve the Blindspot (all 1--100 $p > 0.15$).}
\label{tab:likert100}
\resizebox{\columnwidth}{!}{%
\begin{tabular}{@{}lcccc@{}}
\toprule
Judge & 1--10 $d$ ($p$) & 1--100 $d$ ($p$) & Tie 1--10 & Tie 1--100 \\
\midrule
GPT-4o & --- & $-$.038 (.725) & --- & 65.9\% \\
Qwen3-32B & .012 (.908) & .153 (.154) & 72.7\% & 51.1\% \\
Llama-70B$^a$ & .000 (1.00) & .085 (.440) & 72.7\% & 39.8\% \\
\bottomrule
\multicolumn{5}{l}{\footnotesize $^a$ 5 pairs returned score $= 0$ (out of range); effective $n = 83$.}
\end{tabular}}
\end{table}

\section{NLI Gate Ablation}
\label{app:nli_ablation}

Table~\ref{tab:nli_ablation} compares BT estimates across NLI-pass, NLI-fail, and all pairs.
For formality, NLI-fail pairs reverse to 41.7\% (preferring casual over formal), demonstrating that the gate removes content-confounded pairs rather than amplifying style signal.
The small NLI-fail sample ($n \approx 12$) limits precision; this serves as a bounding analysis.
For verbosity, NLI-fail pairs remain at 75.0\%, consistent with the dominant style effect (OR $= 13.10$) overwhelming content noise.

\begin{table}[h]
\centering
\small
\caption{NLI gate ablation: BT win probability by NLI subset.}
\label{tab:nli_ablation}
\begin{tabular}{lccc}
\toprule
Axis & NLI-pass & NLI-fail & All \\
\midrule
Formality & 66.5\% & 41.7\% & 63.0\% \\
Verbosity & 77.3\% & 75.0\% & 76.0\% \\
\bottomrule
\end{tabular}
\end{table}

\section{Rewriting Artifact Control}
\label{app:artifact_control}

Same-style double-rewrite (DR) controls quantify the rewriting artifact using BT models (consistent with all other pairwise analyses).
Table~\ref{tab:artifact} reports BT $P(\text{original})$, the probability the judge prefers the original over its same-style DR.
Both judges systematically prefer LLM-polished text across axes; for verbosity, BT $P(\text{original}) < 0.17$ indicates that rewriting quality dominates when content is held constant, validating the GEE decomposition that separates style from artifact effects.
Same-style double-rewrite controls for style-neutral polishing artifacts but cannot exclude style-specific polishing interactions (e.g., formal rewriting may polish differently than casual rewriting).

\begin{table}[h]
\centering
\small
\caption{Rewriting artifact: BT $P(\text{orig})$ vs.\ same-style double-rewrite (bootstrap 95\% CI, 10{,}000 resamples). Values ${<}0.5$: judges prefer LLM-polished text.}
\label{tab:artifact}
\begin{tabular}{@{}llcc@{}}
\toprule
Axis & Judge & BT $P(\text{orig})$ [CI] & $p$ \\
\midrule
Form.\ & GPT-4o & .418 [.342, .500] & .037 \\
 & Qwen3 & .377 [.303, .451] & .0006 \\
\midrule
Verb.\ & GPT-4o & .167 [.084, .262] & ${<}.0001$ \\
 & Qwen3 & .082 [.035, .140] & ${<}.0001$ \\
\bottomrule
\end{tabular}
\end{table}

\section{Implementation Details}
\label{app:implementation}

\paragraph{Models and APIs.}
Style-controlled rewriting uses Qwen3-32B (temperature $0.7$, top-$p$ $0.9$).
Target judges: GPT-4o (\texttt{gpt-4o-2024-08-06}, accessed April--May 2026), Qwen3-32B (locally hosted), Llama-3.3-70B (locally hosted via vLLM).
Claude Sonnet 4 (\texttt{claude-sonnet-4-20250514}, accessed May 2026) provides discriminative validity evidence with benchmark-dependent sensitivity.
All judge calls use temperature $0$ and $\texttt{max\_tokens} = 2048$.

\paragraph{NLI gate.}
Content preservation is verified using DeBERTa-v3-large \citep{he2021debertav3} fine-tuned on MNLI, with entailment threshold $\tau = 0.90$.
Bidirectional entailment is applied for formality and register; direction-aware entailment for verbosity (Section~\ref{sec:nli_gate}).

\paragraph{Statistical methods.}
Bradley-Terry parameters are estimated via spectral initialization followed by minorization-maximization \citep{hunter2004mm} using the \texttt{choix} Python library (default L2 regularization $\alpha = 10^{-6}$); varying $\alpha$ from $10^{-6}$ to $1.0$ changes all BT probabilities by ${<}0.002$.
All bootstrap CIs use 10{,}000 pair-level resamples with seed $42$.
GEE models use exchangeable working correlation with robust (sandwich) standard errors via \texttt{statsmodels} (Section~\ref{sec:direction}; diagnostics in Appendix~\ref{app:gee_diagnostics}).

\paragraph{Benchmarks.}
LLMBar \citep{zeng2023llmbar}: 100 instruction-response pairs ($n = 88$ after NLI gate, $n = 355$ for scaled replication).
AlpacaEval \citep{dubois2025lc}: 100 sampled pairs.
MT-Bench \citep{zheng2023judging}: 80 multi-turn pairs.


\section{Holm-Bonferroni Sensitivity Analysis}
\label{app:holm8}

Table~\ref{tab:holm8} reports adjusted $p$-values under the full 8-test family (4 judges $\times$ 2 primary axes), including Gemini-2.5-Flash which is excluded from the primary 6-test family (\S\ref{sec:stat_framework}).
All results remain significant at $\alpha = 0.05$.

\begin{table}[h]
\centering
\small
\caption{Holm-Bonferroni sensitivity analysis: 8-test family (4 judges $\times$ 2 primary axes). All results remain significant under the expanded family ($\alpha = 0.05$).}
\label{tab:holm8}
\resizebox{\columnwidth}{!}{%
\begin{tabular}{@{}cllcccc@{}}
\toprule
Rank & Judge & Axis & BT $P$ & Raw $p$ & Threshold & Adj.\ $p$ \\
\midrule
1 & Gemini Flash & Verb.\ & .708 & ${<}.0001$ & .00625 & ${<}.001$ \\
2 & GPT-4o & Verb.\ & .760 & .0001 & .00714 & .0007 \\
3 & Qwen3-32B & Form.\ & .665 & .0001 & .00833 & .0007 \\
4 & Qwen3-32B & Verb.\ & .773 & .0001 & .0100 & .0007 \\
5 & Llama-70B & Form.\ & .695 & .0001 & .0125 & .0007 \\
6 & Llama-70B & Verb.\ & .735 & .0001 & .0167 & .0007 \\
7 & GPT-4o & Form.\ & .630 & .0004 & .0250 & .0008 \\
8 & Gemini Flash & Form.\ & .574 & .0306 & .0500 & .0306 \\
\bottomrule
\end{tabular}}
\end{table}

\section{Cross-Axis Contamination Analysis}
\label{app:cross_axis_contamination}

Our per-axis design assumes that rewriting along one style axis does not substantially alter the text along other axes.
We quantify cross-axis contamination by measuring off-target surface proxies before and after rewriting.
For formality rewrites ($n = 355$), we measure word count and sentence count as verbosity proxies; for verbosity rewrites ($n = 86$), we measure type-token ratio (TTR) as a lexical diversity proxy for formality, and Flesch-Kincaid (FK) grade level.

\begin{table}[h]
\centering
\small
\caption{Cross-axis contamination: effect of rewriting along one style axis on surface proxies of the other axis. Cohen's $d$: $|d| < 0.20$ negligible, $0.20$--$0.50$ small contamination \citep{cohen1988statistical}.}
\label{tab:cross_axis}
\resizebox{\columnwidth}{!}{%
\begin{tabular}{@{}llrrl@{}}
\toprule
Direction & Proxy & $\Delta$ & $d$ & Interp.\ \\
\midrule
Form.\ $\to$ Verb.\ & Word count & $-3.4\%$ & $-0.23$ & Small \\
Form.\ $\to$ Verb.\ & Sent.\ count & $-0.2\%$ & $-0.03$ & Negligible \\
Verb.\ $\to$ Form.\ & TTR & $-0.006$ & $-0.04$ & Negligible \\
Verb.\ $\to$ Form.\ & FK grade & --- & $0.50$ & Confounded$^*$ \\
\bottomrule
\multicolumn{5}{l}{\footnotesize $^*$ FK $=$ $f$(sentence length, word length); sentence length change is}\\
\multicolumn{5}{l}{\footnotesize inherent to verbosity manipulation, not formality contamination.}
\end{tabular}}
\end{table}

Formality rewrites leave sentence structure virtually unchanged ($d = -0.03$); the word count shift ($d = -0.23$, exceeding the negligible threshold but within the small range) reflects word-level adjustments (contractions, elaboration) that are directionally symmetric (formal$\to$casual $-5.0\%$, casual$\to$formal $+5.6\%$) and do not constitute verbosity contamination.
Verbosity rewrites leave lexical diversity unchanged (TTR $d = -0.04$); the elevated FK grade ($d = 0.50$) is an artifact of FK's dependence on sentence length, which changes by design in verbosity manipulation.
These results confirm that cross-axis contamination is small to negligible for the proxies measured; the GEE originality term (\S\ref{sec:direction}) provides additional protection by absorbing any residual quality confound.
A robustness check adding word-count change as a GEE covariate leaves all Style ORs virtually unchanged ($\Delta < 12\%$, all covariate $p > .28$), confirming that the small cross-axis word-count shift does not confound formality estimates.

\section{FK Length-Control Analysis}
\label{app:fk_length}

The elevated FK grade shift for verbosity rewrites ($d = 0.50$; Table~\ref{tab:cross_axis}) could reflect genuine formality contamination or simply FK's structural dependence on sentence length.
We disentangle these via length-controlled regression and mediation analysis.

\begin{table}[h]
\centering
\small
\caption{Descriptive statistics for FK grade and sentence length changes by rewriting axis. Computed on all pairs; regression analysis below uses NLI-verified pairs only ($n_{\text{form}} = 88$, $n_{\text{verb}} = 86$).}
\label{tab:fk_desc}
\begin{tabular}{@{}llrr@{}}
\toprule
Axis & Metric & $\Delta$ & $d$ \\
\midrule
Verbosity & FK grade & $+2.26$ & $0.50$ \\
Verbosity & Sent.\ length (words) & $+3.40$ & $0.35$ \\
Formality & FK grade & $-0.30$ & $-0.08$ \\
Formality & Sent.\ length (words) & $-0.11$ & $-0.04$ \\
\bottomrule
\end{tabular}
\end{table}

Regressing FK change on style axis without controls yields $\beta = 2.69$ ($p < .001$, $R^2 = 0.100$); adding sentence length as a covariate reduces the coefficient to $\beta = 1.07$ ($p = .020$, $R^2 = 0.546$), a 60.3\% reduction.
Sobel mediation test confirms that sentence length mediates 60\% of the total effect ($z = 2.15$, $p = .031$); the 40\% residual direct effect reflects multi-syllable vocabulary shifts (e.g., ``use'' $\to$ ``utilize'') that are inherent to verbose style rather than formality contamination.

\paragraph{GEE with FK covariate.}
As a further robustness check, we add FK grade change as a covariate to the GEE formality model.
FK grade is not independently significant for any judge ($p > .11$), but its inclusion reduces formality Style ORs by 20--51\% due to collinearity between formality and readability: GPT-4o Style OR drops from $2.70$ to $2.15$ ($p = .132$, no longer significant); Qwen3-32B from $4.20$ to $3.38$ ($p = .014$); Llama-70B from $5.52$ to $2.69$ ($p = .037$).
Two of three judges retain significant formality effects after readability control, confirming that formality style sensitivity is not fully reducible to readability differences.

\section{Human Annotation Agreement Statistics}
\label{app:annotation_agreement}

Two independent annotators evaluated 100 NLI-scored pairs (50 formality, 50 verbosity), stratified into 70 NLI-pass and 30 NLI-fail pairs for content preservation.
Annotator~1 (A1) judged 99/100 pairs content-preserving; Annotator~2 (A2) judged 97/100.

\paragraph{Agreement.}
On the \emph{preserves\_meaning} label, 96/100 pairs are concordant, yielding Gwet's AC$_1 = 0.958$ \citep{gwet2008ac1}.
Cohen's $\kappa = {-}0.015$ is misleading here: with ${>}96\%$ positive prevalence, chance agreement is near-ceiling and $\kappa$ is compressed toward zero regardless of actual agreement, the well-documented prevalence paradox \citep{feinstein1990kappa}.
AC$_1$ is the appropriate statistic under such extreme marginal distributions.

\paragraph{Gate pass precision.}
Both annotators independently confirm 69/70 gate-passed pairs as content-preserving ($98.6\%$; Wilson CI $[92.3\%, 99.8\%]$).
Only two gate-passed pairs were flagged by either annotator (S073 and S063, both formality); excluding these changes BT by ${<}0.3$\,pp and Style OR by ${<}4\%$ relative, with all significance levels unchanged.

\paragraph{Gate-failed analysis.}
Among 30 gate-rejected pairs, A1 judges 30/30 and A2 judges 28/30 as content-preserving, confirming that the gate's dominant error mode is false rejection rather than false acceptance.
This conservative calibration ensures high pass precision at the cost of reduced recall (59.5\%).

\paragraph{Disagreement analysis.}
Four pairs show annotator disagreement: S073, S098, S063, S136 (3 formality, 1 verbosity; 2 NLI-pass, 2 NLI-fail).
All disagreements involve borderline meaning-equivalence judgments rather than clear content contamination.

\paragraph{Comparison with 50-pair pilot.}
The expanded 100-pair annotation tightens the precision estimate from $84$--$100\%$ (individual annotators on 25 gate-passed pairs) to $98.6\%$ (both annotators concordant on 70 gate-passed pairs), and raises AC$_1$ from $0.65$ to $0.96$.

\paragraph{Style detection agreement.}
Agreement on \emph{style\_change\_detected} is lower ($58\%$, $\kappa = 0.138$, AC$_1 = 0.206$), reflecting genuine difficulty in judging whether a rewrite altered style.
This does not affect the primary validation target, which concerns content preservation rather than style detection.

\section{Cross-Benchmark Replication}
\label{app:cross_bench}

Table~\ref{tab:cross_bench} reports cross-benchmark replication results on AlpacaEval and MT-Bench.
Selection criterion: we targeted the six judge$\times$axis cells with the largest LLMBar BT effects (BT $\geq .63$, all $p < .001$), prioritizing cells where the Blindspot was strongest.
The remaining 18/24 cells (3 benchmarks $\times$ 4 target judges $\times$ 2 axes, minus 6 tested) were not evaluated; generalization to weaker-effect configurations remains open.

\begin{table}[h]
\centering
\small
\caption{Cross-benchmark replication on AlpacaEval (AEval) and MT-Bench (MTB). $\ddagger$: both BT and Likert significant; $\dagger$: Blindspot (BT significant, Likert not).}
\label{tab:cross_bench}
\resizebox{\columnwidth}{!}{\begin{tabular}{@{}lllccc@{}}
\toprule
Bench & Judge & Axis & BT $P$ & Likert $d$ ($p$) & Cat.\ \\
\midrule
AEval & GPT-4o & Form.\ & .727$^{***}$ & .005 (.909) & $\dagger$ \\
 & Qwen3 & Form.\ & .657$^{***}$ & .081 (.198) & $\dagger$ \\
 & Qwen3 & Verb.\ & .841$^{***}$ & .375 (.0003) & $\ddagger$ \\
 & Llama-70B & Form.\ & .750$^{***}$ & .212 (.012) & $\ddagger$ \\
\midrule
MTB & GPT-4o & Form.\ & .692$^{***}$ & .111 (.016) & $\ddagger$ \\
 & Qwen3 & Form.\ & .712$^{***}$ & .149 (.007) & $\ddagger$ \\
\bottomrule
\multicolumn{6}{l}{\footnotesize $^{***}$ $p < 0.001$. $\ddagger$: both sig.;\; $\dagger$: Blindspot.}
\end{tabular}}
\end{table}

\section{Claude Cross-Benchmark Analysis}
\label{app:claude_crossbench}

Claude Sonnet~4 exhibits near-null style sensitivity on LLMBar (\S\ref{sec:amplification}; Appendix~\ref{app:claude_tost}).
Cross-benchmark GEE decomposition reveals that Claude's elevated aggregate BT on AEval/MT-Bench reflects quality detection rather than style preference.
Table~\ref{tab:claude_crossbench} reports results across all five tested configurations.

\begin{table}[h]
\centering
\small
\caption{Claude Sonnet~4 cross-benchmark results with GEE decomposition. Formality Style OR is non-significant across all three benchmarks; elevated BT on AEval/MT-Bench is driven by Orig OR (quality detection). AEval verbosity is the sole genuine style effect.}
\label{tab:claude_crossbench}
\resizebox{\columnwidth}{!}{%
\begin{tabular}{@{}llccccc@{}}
\toprule
Bench & Axis & BT $P$ & Likert $d$ ($p$) & Style OR ($p$) & Orig OR ($p$) \\
\midrule
LLMBar & Form.\ & .562 & $-$.021 (.721) & 1.47 (.276) & 6.66 (${<}.001$) \\
LLMBar & Verb.\ & .541 & $-$.039 (.743) & 1.52 (.286) & 0.80 (.563) \\
\midrule
AEval & Form.\ & .698$^{***}$ & .081 (.090) & 1.99 (.082) & 9.33 (${<}.001$) \\
AEval & Verb.\ & .683$^{***}$ & .252 (.002) & \textbf{4.76} (${<}.001$) & 0.87 (.732) \\
\midrule
MTB & Form.\ & .733$^{***}$ & .233 (${<}.001$) & 1.64 (.364) & 11.85 (${<}.001$) \\
\bottomrule
\multicolumn{6}{l}{\footnotesize $^{***}$ $p < 0.001$. All Style ORs except AEval verbosity are non-significant.}\\
\multicolumn{6}{l}{\footnotesize Negative $d$: styled variant scored lower.}
\end{tabular}}
\end{table}

\paragraph{Formality.}
Claude's formality Style OR is non-significant across all three benchmarks ($1.47$/$1.99$/$1.64$, all $p > .08$), indicating consistently low formality sensitivity regardless of benchmark.
The elevated BT on AEval (69.8\%) and MT-Bench (73.3\%) is driven entirely by Orig OR ($9.33$/$11.85$, both $p < .001$): Claude detects quality differences between original and counterfactual texts without exhibiting formality preference.
AEval formality produces a Blindspot (BT significant, Likert n.s.), but GEE decomposition reveals this as quality detection rather than style sensitivity, illustrating why aggregate BT alone is insufficient for bias attribution.

\paragraph{Verbosity.}
AEval verbosity is the sole configuration where Claude exhibits genuine style sensitivity (Style OR $= 4.76$, $p < .001$; Orig OR $= 0.87$, n.s.), a style-dominated pattern.
LLMBar verbosity remains null (Style OR $= 1.52$, n.s.), confirming benchmark-dependent sensitivity.

\section{Rewriting Prompts}
\label{app:rewriting_prompts}

Below are the axis-specific rewriting prompts used in Section~\ref{sec:rewriting}. All prompts are issued to Qwen3-32B-Instruct at temperature 0.7.

\paragraph{Formality.}
\texttt{Rewrite the following text to be more [formal/casual] while preserving the exact same content and meaning. Do not add, remove, or change any factual information. Only adjust the level of formality in the language.}

\paragraph{Verbosity.}
\texttt{Rewrite the following text to be more [verbose/concise] while preserving the exact same content and meaning. Do not add, remove, or change any factual information. Only adjust the level of detail and elaboration.}

\paragraph{Register.}
\texttt{Rewrite the following text to use a more [academic/conversational] register while preserving the exact same content and meaning. Do not add, remove, or change any factual information. Only adjust the register and tone.}

\paragraph{Same-style double-rewrite.}
\texttt{Rewrite the following text in the same style (keep it [formal/casual/verbose/concise]) while preserving the exact same content and meaning. Paraphrase the text without changing its style.}

\section{Non-Tie Concordance Analysis}
\label{app:concordance}

To adjudicate between measurement attenuation and construct divergence as explanations for the Likert Blindspot (\S\ref{sec:likert_pairwise}), we examine whether Likert and pairwise evaluations agree \emph{when Likert does differentiate}.
For each pair where Likert assigned different scores to the two texts (non-tie), we check whether the higher-scoring text matches the pairwise-preferred text; if the two methods measure the same underlying construct, non-tie concordance should substantially exceed chance (50\%).
We additionally compute point-biserial correlations ($r_{\text{pb}}$) between Likert score differences and pairwise outcomes across all pairs with valid pairwise preferences.

\begin{table}[h]
\centering
\small
\caption{Non-tie concordance between Likert and pairwise preferences. For each judge$\times$axis condition, we identify Likert non-tie pairs and check agreement with pairwise preference direction. Wilson 95\% CIs; $p$ from two-sided binomial test against 50\%. $r_{\text{pb}}$: point-biserial correlation between Likert score difference and pairwise outcome (all pairs with valid pairwise preference).}
\label{tab:concordance}
\resizebox{\columnwidth}{!}{%
\begin{tabular}{@{}llrcccrc@{}}
\toprule
Bench & Config & $n_{\text{nt}}$ & Conc. & Rate & 95\% CI & $p$ & $r_{\text{pb}}$ ($p$) \\
\midrule
LLMBar & Llama verb.\ & 25 & 22 & 88.0\% & [70.0, 95.8] & ${<}.001$ & .412 (.001) \\
 & GPT-4o verb.\ & 30 & 25 & 83.3\% & [66.4, 92.7] & ${<}.001$ & .373 (.003) \\
 & Llama form.\ & 18 & 15 & 83.3\% & [60.8, 94.2] & .008 & .447 (.001) \\
 & Qwen3 verb.\ & 27 & 20 & 74.1\% & [55.3, 86.8] & .019 & .296 (.019) \\
 & Gemini verb.\ & 24 & 17 & 70.8\% & [50.8, 85.1] & .064 & .346 (.012) \\
 & Qwen3 form.\ & 15 & 10 & 66.7\% & [41.7, 84.8] & .302 & .274 (.065) \\
 & Gemini form.\ & 15 & 10 & 66.7\% & [41.7, 84.8] & .302 & .373 (.012) \\
\midrule
AEval & Qwen3 form.\ & 26 & 20 & 76.9\% & [57.9, 89.0] & .009 & .203 (.146) \\
 & GPT-4o form.\ & 21 & 15 & 71.4\% & [50.0, 86.2] & .078 & .340 (.019) \\
\midrule
\multicolumn{2}{@{}l}{\textbf{Pooled}} & \textbf{201} & \textbf{154} & \textbf{76.6\%} & \textbf{[70.3, 81.9]} & $\mathbf{{<}.001}$ & --- \\
\bottomrule
\end{tabular}}
\end{table}

Pooled concordance is 76.6\% (154/201 non-tie pairs; binomial $p < .001$; Table~\ref{tab:concordance}), substantially exceeding the 50\% chance baseline.
Point-biserial correlations are positive in all nine conditions ($r_{\text{pb}} = 0.20$--$0.45$), with seven of nine reaching significance ($p < .05$).
Concordance is highest for verbosity ($74$--$88\%$) and for Llama-70B across both axes ($83$--$88\%$), consistent with Llama's lower Likert tie rate (Table~\ref{tab:likert100}).

These results favor the measurement attenuation interpretation: when Likert scoring does differentiate between texts, its preference direction aligns with pairwise judgment in more than three-quarters of cases.
The primary driver of the Blindspot is thus the high Likert tie rate (60--73\% across conditions) rather than a fundamental disagreement between the two evaluation paradigms about what constitutes quality.
The ${\sim}23\%$ discordance rate leaves room for partial construct divergence (pairwise comparison may access quality dimensions that Likert cannot express), but the dominant pattern is signal attenuation through discretization and ceiling effects.


\section*{Use of AI Assistants}
\label{app:ai-assistants}

We used Claude (Anthropic) as an AI coding and writing assistant throughout this work. Specifically:
\begin{itemize}
    \item \textbf{Code:} drafting experiment scripts, data-processing pipelines, and analysis utilities; debugging and refactoring.
    \item \textbf{Writing:} drafting and revising paper sections, polishing prose, and reformatting LaTeX.
    \item \textbf{Literature:} summarizing related papers to inform the related-work section.
\end{itemize}
